% =====================================================================
\section{관련 연구}
\label{sec:related}
% =====================================================================

\paragraph{Per-scene optimization} 3D Gaussian Splatting~\cite{kerbl20233d}은 명시적 3D Gaussian 집합을 미분 가능한 rasterization으로 최적화해 실시간 렌더링 품질의 novel view synthesis를 달성했다. 원 논문은 densification(clone/split), opacity reset, pruning으로 구성된 adaptive density control을 통해 sparse한 초기 point cloud(COLMAP structure-from-motion~\cite{schoenberger2016sfm})에서 시작해 조밀한 표현으로 성장시킨다. 이 densification 메커니즘의 판단 기준은 view-space 위치 gradient가 크면 해당 영역의 표현이 부족하다는 가정에 근거하는데, 이 가정이 관측이 부족한 조건에서도 성립하는지는 검증된 바 없다. 본 논문 \S\ref{sec:starvation}에서 이 메커니즘을 구현 수준에서 직접 계측해 정량화한다.

한편 3DGS가 입력 view 분포에 편향된 형태로 과적합해 분포 밖 시점에서 아티팩트를 낳는다는 보고~\cite{chen2025splatformer}도 있으며, 이는 본 논문의 실패 분석과 상보적이다. FSGS~\cite{zhu2024fsgs}는 이 문제를 겨냥해 proximity-guided Gaussian unpooling과 사전학습된 monocular depth 모델(MiDaS~\cite{ranftl2022midas}) 기반 pseudo-view depth 정규화를 추가한 sparse-view 특화 변형이다.

\paragraph{Feed-forward reconstruction} pixelSplat~\cite{charatan2024pixelsplat}은 이미지 쌍에서 epipolar attention으로 3D Gaussian을 직접 회귀하는 최초의 확장 가능한 feed-forward 방법을 제시했다. MVSplat~\cite{chen2024mvsplat}은 plane-sweep 기반 cost volume(MVSNet~\cite{yao2018mvsnet} 계열)으로 depth를 추정해 Gaussian 중심을 더 정확히 위치시켰다.

DepthSplat~\cite{xu2025depthsplat}은 사전학습된 monocular depth feature를 cost volume에 융합해 더 넓은 view 수 범위에서도 강인한 depth 추정을 달성했다. 이 계열은 학습 시 사용한 view 수 분포(예: MVSplat의 RE10K 체크포인트는 고정 2-view)를 벗어나면 성능이 저하되는 것으로 알려져 있으나, 그 저하가 \emph{정량적으로 얼마나}이며 \emph{어느 지점에서} per-scene optimization이 이를 역전하는지는 체계적으로 다뤄지지 않았다.

\paragraph{경계 지대를 활용하는 연구} 세 연구가 feed-forward와 반복적 정제(iterative refinement)의 경계를 서로 다른 방식으로 넘나든다. Diff3R~\cite{liu2026diff3r}은 미분 가능한 3DGS 최적화 층을 학습 루프에 넣어 ``정제에 유리한 초기값''을 예측하도록 학습하며, 그 과정에서 표준 test-time optimization이 sparse 조건에서 catastrophic overfitting에 빠진다는 점을 문제로 제시한다. ForeSplat~\cite{li2026foresplat}은 정제 후 품질을 목표로 feed-forward 예측 헤드를 메타학습시켜, 예산별 품질 곡선을 개선한다. ReSplat~\cite{xu2026resplat}은 둘과 달리 per-scene gradient 기반 최적화 자체를 쓰지 않는다 — rendering error를 피드백 신호 삼아 recurrent network가 Gaussian을 반복 갱신하도록 학습시켜(4회 반복 후 포화), 단일 forward pass의 성능 한계를 완화하면서도 최적화 기반 3DGS 대비 약 100배 빠른 추론을 유지한다고 보고한다.

세 연구 모두 ``정제(또는 반복 갱신)가 조건에 따라 이득이 달라진다''는 사실을 전제하거나 부수적으로 시사하지만, 그 조건 의존성 자체를 통제된 축(view 수·overlap·예산)에서 체계적으로 매핑하지는 않는다. 실제로 ReSplat은 입력 view 수를 늘리는 부록 실험에서 표본 추출 영역을 함께 확대해 평가 view 자체가 설정마다 달라진다고 스스로 밝히며, 그 결과로도 view 수가 8에서 32로 늘수록 자신과 3DGS의 PSNR 격차가 단조 축소되고(+2.76 → +1.63 → +0.44\,dB) 32-view에서는 LPIPS가 이미 역전됨을 보고한다 — 통제되지 않은 교란 속에서도 본 연구가 다루는 역전 경계의 존재를 시사하는 정황이다. 본 연구는 이 지점을 정면으로, 통제된 조건 아래 다룬다.

\paragraph{Sparse-view 벤치마크와의 차이} 기존 sparse-view 벤치마크는 대개 고정된 view 수(3--12장 내외)에서 단일 시간 예산(대개 완전 수렴)으로 방법들을 비교한다. 본 연구는 view 수, overlap, \emph{계산 시간 예산}을 모두 독립 통제 축으로 다뤄 ``시간을 더 주면 optimization이 feed-forward를 언제 따라잡는가''에 직접 답한다. 또한 overlap을 view 수와 분리된 별도 축으로 명시적으로 통제한다(\S\ref{sec:overlap}).
